\chapter[Band 48: Axiomatische Mengenlehre II]{Band 48 auf einen Blick}
\noindent\textbf{Von den Mengenaxiomen zu Unabhängigkeitsbeweisen}\par\medskip
\noindent\emph{Lesart.}
Dieser Band setzt die Mengenlehre aus Band~3 und die bereits
entwickelten Begriffe der Funktionen, Wohlordnungen und reellen
Zahlen voraus. Er wiederholt diese Grundlagen nicht, sondern
ergänzt die Metatheorie und die beiden Modellkonstruktionen für
die Unabhängigkeitsfragen der Kapitel~4 und~5.

\section*{Kerninhalt}
\begin{itemize}
 \item Syntax, Erfüllung, Korrektheit und ein Beweis des
       Vollständigkeitssatzes durch Henkin-Erweiterungen.
 \item Ordinale, transfinite und fundierte Rekursion,
       Mostowski-Kollaps, Kardinalarithmetik und Skolemhüllen.
 \item Konstruktible Hierarchie, Reflexion, Kondensation und
       relative Konsistenz von ZFC mit GCH.
 \item Mengenforcing über beliebigen abzählbaren Modellen,
       Wahrheitslemma, Erhaltung von ZFC und ccc-Kardinalerhaltung.
 \item Cohen-Forcing und Zählung schöner Namen bis zum Modell
       mit \(2^{\aleph_0}=2^{\aleph_1}=\aleph_2\).
\end{itemize}

\section*{Erreichte Ergebnisse}
Die Vorwärtsrichtung
\(A\approx B\Rightarrow\mathcal P(A)\approx\mathcal P(B)\),
Cantors strikte Ungleichung und die endliche Rückrichtung
werden in Beweistabellen hergeleitet.

Unter der ausdrücklich genannten Voraussetzung
\(\operatorname{Con}(\mathsf{ZFC})\) ist die universelle Aussage
\[
 \forall A,B\
 (\mathcal P(A)\approx\mathcal P(B)\to A\approx B)
\]
weder aus ZFC beweisbar noch in ZFC widerlegbar; siehe
\FormulaRefAuto{B48PCIndependence}[thm].
Dasselbe gilt für die Kontinuumshypothese
\(2^{\aleph_0}=\aleph_1\); siehe
\FormulaRefAuto{B48CHIndependence}[thm].
Die Gleichheit \(|\mathbb R|=2^{\aleph_0}\) und die
übliche Zwischenmächtigkeitsformulierung werden ebenfalls bewiesen.

\section*{Didaktische Stellung}
Der Aufbau folgt dem Schema Definition, Hilfssatz,
Beweistabelle und Anwendung. Der Unvollständigkeitssatz ist
keine Voraussetzung. Weder die Konsistenz von ZFC noch
die Existenz eines transitiven ZFC-Modells wird
voraussetzungslos behauptet. Die Tabellen sind mathematische
Beweise, keine maschinell geprüfte Formalisierung.

